\section{Black-Scholes, Volatilité Implicite et Grecques}
\label{sec:bs_iv}

\subsection{Formule de Black 1976}

Pour les options sur futures, on utilise la formule de Black (1976) :
\begin{align}
  C &= e^{-rT}\bigl[F_T\,\N(d_1) - K\,\N(d_2)\bigr] \label{eq:black_call}\\
  P &= e^{-rT}\bigl[K\,\N(-d_2) - F_T\,\N(-d_1)\bigr] \label{eq:black_put}
\end{align}
avec $d_1 = \dfrac{\ln(F_T/K) + \frac{1}{2}\sigma^2 T}{\sigma\sqrt{T}}$ et
$d_2 = d_1 - \sigma\sqrt{T}$.

\subsection{Extraction de la Volatilité Implicite}

La volatilité implicite $\iv$ est la valeur de $\sigma$ telle que
$\text{Black}(F_T, K, T, r, \sigma) = m$ où $m$ est le prix de marché.

\paragraph{Newton-Raphson.} Partant de $\sigma^{(0)}$ (approximation de
Brenner-Subrahmanyam), on itère :
\begin{equation}
  \sigma^{(n+1)} = \sigma^{(n)} - \frac{\text{Black}(\sigma^{(n)}) - m}{\text{Vega}(\sigma^{(n)})}
\end{equation}
Convergence quadratique (~5 itérations). Fallback dichotomie si Vega $< \varepsilon$
ou divergence.

\subsection{Grecques}

\begin{align}
  \Delta_C &= e^{-rT}\N(d_1) & \Delta_P &= -e^{-rT}\N(-d_1)\\
  \Gamma   &= \frac{e^{-rT}\,n(d_1)}{F_T\,\sigma\,\sqrt{T}} &
  \mathcal{V} &= e^{-rT}\,F_T\,n(d_1)\,\sqrt{T}
\end{align}

% TODO : statistiques de convergence Newton vs Dichotomie
% TODO : smile de volatilité pour plusieurs maturités (plot_iv_smile)
